\section{Recovery and output validation audit}
The four original bounded runs were blocked before inference by a missing
\texttt{OLLAMA\_BASE\_URL} environment variable. Each had zero saved or failed
predictions; pending method outputs were 60 for TAT-QA, 60 for TAT-HQA, 60 for
MultiHiertt, and 66 for ConvFinQA. No inference ledger existed for those runs.
The recorded local Windows endpoint was restored and its model digest verified;
a real TAT-QA request completed before bulk dispatch. That exact request is
reused from its ledger, including its original response and attempt metadata.
The frozen IDs, ordering, prompts, evidence cap, schema and decoding settings
were retained. No FinQA generation was repeated.

The saved-output audit examined all 1344 FinQA development rows labelled
``Malformed structured output'', plus 92 partial-test rows, 35 pilot rows and
8 smoke rows. All were valid JSON satisfying the broad response schema, but
violated the kind-specific plan contract. For example, table aggregation
omitted its required row, and extraction returned arithmetic steps instead of
its required operand. Every malformed row failed these declared plan rules;
none passed the rules and was then rejected. The schema admits combinations
that the interpreter disallows; this is a limitation of generation constraints,
not evidence that these returned plans have valid executable meanings.
Relaxing the executor or deleting fields would reinterpret model answers.
Accordingly the original FinQA scores remain unchanged, and no corrected
inference rerun is mixed with them. The generic error label obscures this
distinction; the detailed audit is saved separately.
